\section*{Hoofdstuk \ref{ch:b1:periodicity} --- Periodiciteit}

\begin{solution}{exo:b1:periodicity:1}
\ce{Ca} $[\ce{Ar}]\,4\text{s}^2$: \omterm{def:b1:periodicity:table}{periode} 4, \omterm{def:b1:periodicity:table}{groep} 2, \omterm{def:b1:periodicity:table}{s-blok}. \ce{Fe}
$[\ce{Ar}]\,3\text{d}^6 4\text{s}^2$: \omterm{def:b1:periodicity:table}{periode} 4, \omterm{def:b1:periodicity:table}{groep} 8, \omterm{def:b1:periodicity:table}{d-blok}. \ce{Se}
$[\ce{Ar}]\,3\text{d}^{10} 4\text{s}^2 4\text{p}^4$: \omterm{def:b1:periodicity:table}{periode} 4, \omterm{def:b1:periodicity:table}{groep} 16,
\omterm{def:b1:periodicity:table}{p-blok}. \ce{I} $[\ce{Kr}]\,4\text{d}^{10} 5\text{s}^2 5\text{p}^5$: \omterm{def:b1:periodicity:table}{periode} 5,
\omterm{def:b1:periodicity:table}{groep} 17, \omterm{def:b1:periodicity:table}{p-blok}.
\end{solution}

\begin{solution}{exo:b1:periodicity:2}
\ce{Na} is groter dan \ce{Cl}: zelfde \omterm{def:b1:periodicity:table}{periode}, kleinere $Z^{*}$. \ce{Na}
is groter dan \ce{Na+}, dat zijn 3s-schil kwijt is. \ce{Cl-}
(\qty{181}{pm}) is groter dan \ce{F-} (\qty{133}{pm}): één schil meer.
\ce{Na+} (\qty{102}{pm}) is groter dan \ce{Mg^{2+}} (\qty{72}{pm}):
dezelfde tien elektronen, kernlading 11 tegen 12.
% ledger: rion:Cl-VI, rion:F-VI, rion:Na+VI, rion:Mg2+VI
\end{solution}

\begin{solution}{exo:b1:periodicity:3}
$\ce{Na} (5.14) < \ce{Al} (5.99) < \ce{Mg} (7.65) < \ce{S} (10.36) <
\ce{P} (10.49) < \ce{Ar} (15.76)$. De algemene toename langs de \omterm{def:b1:periodicity:table}{periode}
wordt twee keer doorbroken: \ce{Al} onder \ce{Mg} (het verwijderde
elektron is het eerste 3p-elektron, hoger in energie dan 3s) en \ce{S}
onder \ce{P} (het verwijderde elektron is het eerste gepaarde
3p-elektron).
\end{solution}

\begin{solution}{exo:b1:periodicity:4}
De \omterm{def:b1:periodicity:polarisability}{polariseerbaarheid} meet hoe gemakkelijk de elektronenwolk door een
elektrisch veld wordt vervormd. \ce{I-} is polariseerbaarder dan
\ce{F-}: het is groter en zijn buitenste elektronen zitten verder van de
kern. \ce{Cl-} (18 elektronen vastgehouden door $Z = 17$) is veel
polariseerbaarder dan \ce{Na+} (10 elektronen vastgehouden door
$Z = 11$), zoals anionen dat in het algemeen zijn ten opzichte van
kationen.
\end{solution}

\begin{solution}{exo:b1:periodicity:5}
Verhoudingen: $18.83/5.99 = 3.1$, $28.45/18.83 = 1.5$, $119.99/28.45 = 4.2$.
De sprong komt na het derde elektron: drie \omterm{def:b1:quantum-numbers:configuration}{valentie-elektronen}, \omterm{def:b1:periodicity:table}{groep}
13, \omterm{def:b1:periodicity:table}{periode} 3: aluminium. Het vormt \ce{Al^{3+}}, $1\text{s}^2 2\text{s}^2
2\text{p}^6$.
% ledger: ie:Al, ie2:Al, ie3:Al, ie4:Al
\end{solution}

\begin{solution}{exo:b1:periodicity:6}
Het toegevoegde elektron van \ce{F-} komt in de kleine 2p-schil, waar
het sterk wordt afgestoten door de zeven \omterm{def:b1:quantum-numbers:configuration}{valentie-elektronen} die er al
zitten; in de grotere 3p-schil van chloor is de afstoting kleiner, zodat
er meer energie vrijkomt. In stikstof ($2\text{p}^3$, \omorbs{u,u,u})
zou het extra elektron moeten paren in een half gevulde \omterm{def:b1:quantum-numbers:orbital}{subschil}; in
magnesium ($3\text{s}^2$) zou het de hoger gelegen 3p-subschil moeten
openen. In geen van beide gevallen is het anion stabieler dan het atoom
met een vrij elektron.
\end{solution}

\begin{solution}{exo:b1:periodicity:7}
$\chi_M(\ce{Na}) = (5.14 + 0.55)/2 = \qty{2.85}{eV}$; $\chi_M(\ce{Cl}) =
(12.97 + 3.61)/2 = \qty{8.29}{eV}$. Het verschil is groot: het
elektronenpaar van een \ce{Na-Cl}-binding hoort bijna volledig bij
chloor, en natriumchloride is ionair, \ce{Na+ Cl-}.
\end{solution}

\begin{solution}{exo:b1:periodicity:8}
\ce{H-F}: $\Delta\chi = 1.78$, op de grens van het ionaire gebied (in
de praktijk een zeer polaire covalente binding), \ce{F^{\delta-}}.
\ce{C-H}: 0.35, vrijwel apolair. \ce{Na-Cl}: 2.23, ionair, \ce{Cl^{-}}.
\ce{C-Cl}: 0.61, polair covalent, \ce{C^{\delta+}-Cl^{\delta-}}.
\ce{O-H}: 1.24, polair covalent, \ce{O^{\delta-}-H^{\delta+}}.
\end{solution}

\begin{solution}{exo:b1:periodicity:9}
In de groep naar beneden zit het valentie-elektron in een schil met
grotere $n$, verder van de kern en afgeschermd door meer \omterm{def:b1:quantum-numbers:configuration}{rompelektronen};
het wordt gemakkelijker verwijderd. Rubidium, één schil verder, heeft een
\omterm{def:b1:periodicity:ionisation-energy}{eerste ionisatie-energie} onder \qty{4.34}{eV}.
\end{solution}

\begin{solution}{exo:b1:periodicity:10}
Zink is $[\ce{Ar}]\,3\text{d}^{10} 4\text{s}^2$: zijn elektron komt uit
een volle 4s; gallium, $\dots 4\text{s}^2 4\text{p}^1$, verliest zijn
enige 4p-elektron, dat hoger in energie ligt en afgeschermd wordt door de
volle 4s en 3d. Arseen, $4\text{p}^3$, verliest een elektron uit een half
gevulde \omterm{def:b1:quantum-numbers:orbital}{subschil}; seleen, $4\text{p}^4$, verliest het gepaarde elektron,
dat door de afstoting van zijn partner omhoog wordt geduwd: zijn
ionisatie-energie is iets lager.
\end{solution}

\begin{solution}{exo:b1:periodicity:11}
$\Delta E = (3.98 - 2.20)^2 = 1.78^2 = \qty{3.17}{eV} = 3.17 \times 96.5 =
\qty{306}{kJ/mol}$. Dan is $D(\ce{H-F}) = \tfrac12(436 + 158) + 306 = 297 +
306 \approx \qty{603}{kJ/mol}$. Het grote overschot is de extra
aantrekking tussen de partiële ladingen $\ce{H^{\delta+}}$ en
$\ce{F^{\delta-}}$, het gevolg van het grote verschil in
\omterm{def:b1:periodicity:electronegativity}{elektronegativiteit}.
\end{solution}

\begin{solution}{exo:b1:periodicity:12}
De \omterm{def:b1:periodicity:electronegativity}{schaal van Pauling} is gedefinieerd uit de energieën van bindingen
A--B: een element dat geen binding vormt, heeft geen waarde volgens
Pauling. Helium, neon en argon vormen geen stabiele verbindingen.
Krypton en vooral xenon hebben grotere, polariseerbaardere
valentieschillen en lagere ionisatie-energieën (\qty{14.0}{eV} voor
krypton tegen \qty{15.76}{eV} voor argon); ze worden door fluor en
zuurstof geoxideerd en vormen verbindingen zoals \ce{XeF2}, zodat hun
bindingsenergieën gemeten kunnen worden.
% ledger: ie:Kr, ie:Ar
\end{solution}

\begin{solution}{pb:b1:periodicity:1}
\textbf{1.} Elk elektron wordt verwijderd uit een ion dat positiever en
kleiner is dan het vorige en zijn elektronen dus steviger vasthoudt.
\textbf{2.} $15.04/7.65 = 1.97$; $80.14/15.04 = 5.33$; $109.27/80.14 = 1.36$.
\textbf{3.} De sprong komt na twee elektronen: X heeft twee
\omterm{def:b1:quantum-numbers:configuration}{valentie-elektronen}, \omterm{def:b1:periodicity:table}{groep} 2; in \omterm{def:b1:periodicity:table}{periode} 3 is dat magnesium.
\textbf{4.} \ce{Mg^{2+}}, $1\text{s}^2 2\text{s}^2 2\text{p}^6$, de
configuratie van neon.
\textbf{5.} Een metaal: \omterm{def:b1:periodicity:table}{groep} 2, links in het systeem, met lage eerste
ionisatie-energieën.
\textbf{6.} Het derde elektron komt uit de 2p-romp, met kleinere $n$,
veel dichter bij de kern en nauwelijks afgeschermd, en wordt verwijderd
uit een tweevoudig geladen ion.
% ledger: ie:Mg, ie2:Mg, ie3:Mg, ie4:Mg
\textbf{7.} Alle hebben tien elektronen: $1\text{s}^2 2\text{s}^2 2\text{p}^6$.
\textbf{8.} Dezelfde tien elektronen worden vastgehouden door kernladingen
8, 9, 11, 12, 13: de wolk trekt samen naarmate $Z$ groeit.
\textbf{9.} \ce{O^{2-}}: het grootste, met de laagste kernlading voor zijn
tien elektronen, en een anion.
\textbf{10.} $140/54 = 2.6$.
\textbf{11.} $198.8/2 = \qty{99.4}{pm}$; het anion is $181/99.4 = 1.8$
keer zo groot als de \omterm{def:b1:periodicity:radius}{covalente straal}.
% ledger: re:Cl2, rion:Cl-VI
\textbf{12.} $\ce{Mg^{2+}} < \ce{Na+} < \ce{Cl-}$: de twee kationen zijn
iso-elektronisch (lading 12 tegen 11), en \ce{Cl-} heeft een derde schil.
\textbf{13.} $\chi_M$ = 10.41 (\ce{F}), 8.29 (\ce{Cl}), 7.59 (\ce{Br}),
7.53 (\ce{O}), 6.26 (\ce{C}), in eV.
\textbf{14.} Mulliken: $\ce{F} > \ce{Cl} > \ce{Br} > \ce{O} > \ce{C}$;
Pauling: $\ce{F} > \ce{O} > \ce{Cl} > \ce{Br} > \ce{C}$. Zuurstof schuift
op van de vierde naar de tweede plaats.
\textbf{15.} $3.98/10.41 = 0.382$, $3.16/8.29 = 0.381$, $2.96/7.59 =
0.390$: vrijwel constant, ongeveer 0.38; voor de halogenen zijn de twee
schalen evenredig.
\textbf{16.} De \omterm{def:b1:periodicity:electron-affinity}{elektronenaffiniteit} van zuurstof is klein
(\qty{1.44}{eV}) omdat het toegevoegde elektron moet paren in de compacte
2p-schil; de waarde volgens Mulliken van het vrije atoom onderschat de
trekkracht die zuurstof in zijn bindingen uitoefent, en die meet de
\omterm{def:b1:periodicity:electronegativity}{schaal van Pauling}, opgebouwd uit bindingen, rechtstreeks.
\textbf{17.} Chloor in \ce{HCl}; fluor in \ce{ClF}.
\textbf{18.} $\Delta\chi = 0.96$: polair covalent.
\textbf{19.} $D(\ce{H-H}) = 2 \times 218.0 = \qty{436.0}{kJ/mol}$.
\textbf{20.} $D(\ce{Cl-Cl}) = 2 \times 121.3 = \qty{242.6}{kJ/mol}$.
\textbf{21.} $D(\ce{H-Cl}) = 218.0 + 121.3 - (-92.3) = \qty{431.6}{kJ/mol}$.
\textbf{22.} $\Delta E = 431.6 - \tfrac12(436.0 + 242.6) = 431.6 - 339.3 =
\qty{92.3}{kJ/mol}$. In symbolen: $D(\ce{H-Cl}) = \Delta_fH(\ce{H}) +
\Delta_fH(\ce{Cl}) - \Delta_fH(\ce{HCl})$, terwijl het gemiddelde van de
symmetrische bindingen $\Delta_fH(\ce{H}) + \Delta_fH(\ce{Cl})$ is, zodat
$\Delta E =
-\Delta_fH(\ce{HCl})$.
\textbf{23.} $92.3/96.485 = \qty{0.957}{eV}$.
\textbf{24.} $|\chi(\ce{Cl}) - \chi(\ce{H})| = \sqrt{0.957} =
\textbf{0.98}$, tegen $3.16 - 2.20 = 0.96$ in de tabellen: het rekenwerk
van Pauling, uitgevoerd met moderne gegevens, levert het getabelleerde
verschil op 0.02 na.
\end{solution}
